\documentclass[10pt,a4paper]{amsart}
\usepackage[margin=19mm]{geometry}
\usepackage{amsmath,amssymb,mathtools}
\usepackage[scr=rsfs,cal=euler]{mathalfa}
\usepackage{xfrac}
\usepackage[unicode,hidelinks,bookmarks=false]{hyperref}
\newcommand{\ie}{i.e.\ }
\newcommand{\cf}{cf.\ }
\newcommand{\resp}{respectively\ }
\newcommand{\Subsection}{\subsection}
\providecommand{\nobreakdash}{\nobreak\hbox{-}}
\setlength{\emergencystretch}{2em}
\multlinegap0pt
\usepackage{bm}
\usepackage{amssymb}
\usepackage{mathtools}
\usepackage{breqn}
\usepackage[normalem]{ulem}
\usepackage{dsfont}
\usepackage{float}
\usepackage{xcolor}
\usepackage{tikz}
\usepackage[shortlabels]{enumitem}
\usepackage{algorithm}\usepackage{algpseudocode}
\newtheorem{definition}{Definition}
\newtheorem{lemma}{Lemma}
\usepackage{cleveref}
\crefname{definition}{Definition}{Definitions}
\crefname{lemma}{Lemma}{Lemmas}
\newcommand{\K}{\mathcal K}
\newcommand{\Lip}{\textup{Lip}}
\newcommand{\Om}{\Omega}
\newcommand{\dd}{\mathrm{d}}
\newcommand{\de}{\mathbf{D}}
\newcommand{\pp}{a.e.\,}
\title{A coupled prediction-correction Hughes' model\\ for congested crowd motion}
\author{Jos\'{e} Miguel Urbano$^{\diamond}$}
\date{Source: arXiv:2606.02899v1 (CC BY 4.0)}
\begin{document}
\maketitle

\noindent\emph{Source and licence.} A coupled prediction-correction Hughes' model\\ for congested crowd motion. Version arXiv:2606.02899v1 (CC BY 4.0), distributed by arXiv under the Creative Commons Attribution 4.0 International licence, http://creativecommons.org/licenses/by/4.0/, as recorded in the arXiv licence metadata for this exact version and shown on the abstract page https://arxiv.org/abs/2606.02899v1. Full licence notice: \texttt{licenses/arxiv-2606.02899-CC-BY-4.0.txt}.\par\smallskip

\noindent\textit{Editorial scope.} Counted unit: the article's lemma lem:rho (source0.tex lines 2091--2093). The article's complete original proof of it is delivered with it (source0.tex lines 2095--2133, one \texttt{proof} environment of 39 lines). No mathematics is written, repaired or re-derived: the statement and the proof are the article's own lines, in the article's own notation, and no retained line is substituted or re-flowed.  Retained uncounted as the source-local context the proof needs: lines 2030--2040; lines 2043--2048.  Nothing was omitted from any retained span, so the package carries no omission record.  No retained line contains a citation, so the wrapper carries no bibliography and no bibliography entry.  No retained line contains a figure environment, an image-inclusion command or drawing code, so no image is delivered and the package declares no asset dependency.  Editorial formatting only, in addition: an AMS \texttt{amsart} preamble that restates the article's own declarations and macros used by the retained lines, namely the environments \texttt{definition}, \texttt{lemma} on separate counters, together with the standard packages those lines need.  The retained environments print in this document as Definition 1, Lemma 1 in order of appearance, whereas the article prints the same items with the numbering of its own full document.  The complete native source of the article as distributed in the arXiv e-print for this exact version is bundled unchanged as \texttt{sources/201939/source0.tex}.
\begin{definition}\label{def:interpolants}
Given $\epsilon >0$, we define, for $t \in (t_i, t_{i+1}]$, and $i = 0, \dots, n-1,$
\begin{equation}\label{eq:approx-sol}
	    \rho^\epsilon(t) = \rho^{i+1}, \quad p^\epsilon(t) = p^{i+1}, \quad \text{and} \quad \de^\epsilon(t) = \de^{i+1},
\end{equation}
with the initial data defined at $t=0$ as $\rho^\epsilon(0) = \rho^0$, $p^\epsilon(0) = p^0$, and $\de^\epsilon(0) = \de^0$. 
Moreover, we define the piecewise linear interpolant function $\tilde{\rho}^\epsilon$ as
\begin{equation}\label{eq:tilde-rho}
    \tilde{\rho}^\epsilon(t) = \frac{t - t_i}{\epsilon}\rho^{i+1} + \frac{t_{i+1} - t}{\epsilon}\rho^{i}, \quad \text{for } t \in (t_i, t_{i+1}].
\end{equation}
\end{definition}

\begin{lemma}\label{lem:d-ineq}
	For all $\epsilon>0$, any test function $\xi\in\K$ and any $k>0$, the approximate solutions defined in \cref{def:interpolants} satisfy the following inequality for \pp $t\in (0,T)$:
	\begin{equation}\label{eq:final-discrete-ineq}
    \int_{\Omega} \partial_t \tilde{\rho}^{\epsilon} T_k(\beta^{-1}(\tilde{\rho}^{\epsilon})-\xi) \dd x + \int_{\Omega} \rho^{\epsilon} \nabla \de^{\epsilon} \cdot \nabla T_k(p^{\epsilon}-\xi) \dd x \le \int_{\Omega} f^{\epsilon} T_k(p^{\epsilon}-\xi) \dd x.
\end{equation}
\end{lemma}

\begin{lemma}\label{lem:rho}
    The sequences $(p^\epsilon)_{\epsilon}$ and $(\tilde{\rho}^{\epsilon})_{\epsilon}$ are bounded in $L^{\infty}(0,T; \Lip_{D}(\Omega))$, and the sequence $(\partial_t \tilde{\rho}^{\epsilon})_{\epsilon}$ is bounded in $L^1(0,T; \Lip_{D}^*(\Om))$.
\end{lemma}

\begin{proof}
First, we note that $p^{\epsilon}(t) \in \K$, for all $t \in [0,T]$. Since the space $\K$ is a bounded subset of $\Lip_D(\Omega)$ and $L^\infty(\Omega)$, the sequence $(p^{\epsilon})_{\epsilon}$ is bounded in $L^{\infty}(0,T; \Lip_D(\Omega))$. Moreover, since $\rho^\epsilon = \beta(p^{\epsilon})$ \pp and $\beta$ is Lipschitz continuous, we deduce that $(\rho^\epsilon)_{\epsilon}$ is bounded in $L^{\infty}(0,T; \Lip_{D}(\Omega))$. By convex combination, the piecewise linear interpolant $(\tilde{\rho}^{\epsilon})_{\epsilon}$ is itself bounded in $L^\infty(0,T; \Lip_{D}(\Omega))$.

     To bound the time derivative $\partial_t \tilde{\rho}^{\epsilon}$, we observe that since $\K$ is symmetric, $-\xi \in \K$ for any $\xi \in \K$. Moreover, going back to the proof of \cref{lem:d-ineq}, we have  for all $t \in (t_i, t_{i+1}]$,
\begin{equation}\label{eq:v_ineq_lemrho}
    \int_{\Omega} \partial_t \tilde{\rho}^{\epsilon} T_k (p^\epsilon-\xi) \dd x + \int_{\Omega} \rho^\epsilon \nabla \de^\epsilon \cdot \nabla T_k(p^\epsilon-\xi) \dd x \le \int_{\Omega} f^\epsilon T_k(p^\epsilon-\xi) \dd x, \quad \text{for all } \xi \in \K.
\end{equation}
    Since $p^\epsilon$ and $\xi$ belong to $\K \subset L^\infty(\Omega)$, taking $k > \Vert p^\epsilon - \xi \Vert_\infty$ ensures that $T_k(p^\epsilon - \xi) = p^\epsilon - \xi$. Thus, the inequality \eqref{eq:v_ineq_lemrho} simplifies to
\begin{equation}\label{eq:v_ineq_k_large}
    \int_{\Omega} \partial_t \tilde{\rho}^{\epsilon} (p^\epsilon-\xi) \dd x + \int_{\Omega} \rho^{\epsilon}\nabla \de^\epsilon \cdot \nabla(p^{\epsilon}-\xi) \dd x \le \int_{\Omega} f^{\epsilon} (p^{\epsilon}-\xi) \dd x,
\end{equation}
     for any $\xi \in \K$. Substituting $-\xi$ for $\xi$ in the \eqref{eq:v_ineq_k_large}, we obtain
    \begin{equation}
        \int_{\Omega} \partial_t \tilde{\rho}^{\epsilon} (\beta^{-1}(\tilde{\rho}^{\epsilon})+\xi) \dd x + \int_{\Omega} \rho^{\epsilon}\nabla \de^\epsilon \cdot \nabla(p^{\epsilon}+\xi) \dd x \le \int_{\Omega} f^{\epsilon} (p^{\epsilon}+\xi) \dd x.
    \end{equation}	
    Expanding the first term yields
    \begin{equation}\label{eq:bound_xi}
        \int_{\Omega} \partial_t \tilde{\rho}^{\epsilon} \xi \dd x \leq - \int_{\Omega} \rho^{\epsilon}\nabla \de^\epsilon \cdot \nabla(p^{\epsilon}+\xi) \dd x + \int_{\Omega} f^{\epsilon} (p^{\epsilon}+\xi) \dd x - \int_{\Omega} \partial_t \tilde{\rho}^{\epsilon} \beta^{-1}(\tilde{\rho}^{\epsilon}) \dd x, 
    \end{equation}
    for all  $t \in [t_i, t_{i+1})$.
    Since $\partial_t \tilde{\rho}^{\epsilon} \beta^{-1}(\tilde{\rho}^{\epsilon}) = \frac{\dd}{\dd t} \int_{0}^{\tilde{\rho}^{\epsilon}(t)} \beta^{-1}(s) \dd s$, we get, by substituting this into \eqref{eq:bound_xi},
    \begin{equation}\label{eq:eqxx}
        \int_{\Omega} \partial_t \tilde{\rho}^{\epsilon} \xi \dd x \leq -\int_{\Omega} \rho^\epsilon \nabla \de^\epsilon \cdot \nabla(p^\epsilon +\xi) \dd x + \int_{\Omega} f^{\epsilon} (p^{\epsilon}+\xi) \dd x - \frac{\dd}{\dd t} \int_{\Omega} \int_{0}^{\tilde{\rho}^{\epsilon}(t, x)} \beta^{-1}(s) \dd s \dd x, 
    \end{equation}
    for all $t \in [t_i, t_{i+1})$.
    Integrating \eqref{eq:eqxx} over the entire time interval $[0,T]$, we obtain
    \begin{equation*}
        \begin{aligned}
            \int_{0}^{T}\int_{\Omega} \partial_t \tilde{\rho}^{\epsilon} \xi \dd x \dd t 
            &\leq -\int_{0}^{T} \int_{\Omega} \rho^\epsilon \nabla \de^\epsilon \cdot \nabla(p^\epsilon +\xi) \dd x \dd t + \int_{0}^{T}\int_{\Omega} f^{\epsilon} (p^{\epsilon}+\xi) \dd x \dd t \\
            &\quad + \int_{\Omega}\int_{0}^{\rho_0(x)} \beta^{-1}(s)\dd s \dd x - \int_{\Omega} \int_{0}^{\tilde{\rho}^\epsilon(T, x)} \beta^{-1}(s) \dd s \dd x.
        \end{aligned}
    \end{equation*}
    Since $\rho^\epsilon$ and $p^\epsilon$ are bounded in $L^{\infty}(0,T; \Lip_{D}(\Omega))$, $\nabla \de^\epsilon$ is bounded, and $\beta^{-1}$ is continuous, the right-hand side is bounded by a constant $C > 0$ independent of $\epsilon$. Taking the supremum over all test functions $\xi \in \K$, we deduce that
    \begin{equation}
        \sup_{\xi \in \K} \int_{0}^{T} \int_{\Omega} \partial_t \tilde{\rho}^{\epsilon} \xi \dd x \dd t \leq C < \infty,
    \end{equation}
    which means that the sequence $(\partial_t \tilde{\rho}^{\epsilon})_{\epsilon}$ is bounded in $L^1(0,T; \Lip_{D}^*(\Om))$.
\end{proof}

\end{document}
